% 第 5 题三视图（半球＋四棱锥）：按固定源 PDF 矢量路径换算（1 pt → 0.03515 cm）。
% 原点取图组左下（源页 (156.82, 781.34) pt），y 向上。1 单位 = 34.00 pt：三角形底=高=34.00（1 单位）。
% 正（主）视图＝半圆（**平边弦在上、弧向下鼓**，r 24.05＝√2/2 单位，弧底与三角形顶点相接）
% ＋等腰三角形；侧（左）视图与之全等（右移 85.01）；俯视图＝实线圆（r 24.05）＋虚线内接正方形
% （边 34.00）及其两条对角线。尺寸「1」共 3 处（正视图纵、横与侧视图横），俯视图无标注。
% 图注照录源内写法：「正（主）视图」「侧（左）视图」（带括注）、「俯视图」（无括注）。
% 弦（y 619.90）与弧底（643.95）均已由 400 dpi 墨迹带实测复核（弦线 48.6 pt 连续实线）。
\begin{tikzpicture}[>=stealth,every node/.style={font=\small},x=0.03515cm,y=0.03515cm]
% ===== 正（主）视图 =====
\draw (0,161.44)--(48.09,161.44);
\draw (48.09,161.44) arc[start angle=0,end angle=-180,radius=24.05];
\draw (24.04,137.39)--(7.04,103.39)--(41.04,103.39)--cycle;
% 纵向尺寸线（x=1.94，两段让位给标注；端部实心箭头；界外延长线两条）
\draw[->] (1.94,127.19)--(1.94,137.39);
\draw[->] (1.94,113.59)--(1.94,103.39);
\node at (2.0,119.5) {\(1\)};
\draw (-1.46,137.39)--(24.04,137.39);
\draw (-1.46,103.39)--(5.34,103.39);
% 横向尺寸线（y=98.29，两段让位；端部箭头；竖直端线两条）
\draw[->] (18.94,98.29)--(7.04,98.29);
\draw[->] (29.14,98.29)--(41.04,98.29);
\node at (24.12,98.29) {\(1\)};
\draw (7.04,94.89)--(7.04,101.69);
\draw (41.04,94.89)--(41.04,101.69);
% ===== 侧（左）视图（与正视图全等，右移 85.01）=====
\draw (85.01,161.44)--(133.10,161.44);
\draw (133.10,161.44) arc[start angle=0,end angle=-180,radius=24.05];
\draw (109.05,137.39)--(92.05,103.39)--(126.05,103.39)--cycle;
\draw[->] (103.95,98.29)--(92.05,98.29);
\draw[->] (114.15,98.29)--(126.05,98.29);
\node at (109.13,98.29) {\(1\)};
\draw (92.05,94.89)--(92.05,101.69);
\draw (126.05,94.89)--(126.05,101.69);
% ===== 俯视图 =====
\draw (24.04,38.78) circle (24.05);
\draw[dash pattern=on 2.99pt off 2.99pt] (7.04,55.79)--(41.04,55.79)--(41.04,21.78)--(7.04,21.78)--cycle;
\draw[dash pattern=on 2.99pt off 2.99pt] (7.04,21.78)--(41.04,55.79);
\draw[dash pattern=on 2.99pt off 2.99pt] (41.04,21.78)--(7.04,55.79);
% ===== 图注（居中于各子图下方）=====
\node[anchor=base,inner sep=0] at (24.11,73.68) {正（主）视图};
\node[anchor=base,inner sep=0] at (109.15,73.65) {侧（左）视图};
\node[anchor=base,inner sep=0] at (24.11,0) {俯视图};
\end{tikzpicture}
